\documentclass[a4paper,11pt]{article}
\usepackage[nosumlimits]{amsmath}
\usepackage{amssymb}
\usepackage{amsthm}
\usepackage{enumitem}

\setlist[enumerate,1]{label = \emph{\alph*}),
  ref = \theenumi.\emph{\alph*}}
\setlist[enumerate,2]{label = \arabic*.,
  ref = \arabic*}
\setlist[enumerate,3]{label = \roman*),
  ref = \theenumii.\roman*}

\DeclareMathOperator{\dom}{dom}
\DeclareMathOperator{\ran}{ran}
\DeclareMathOperator{\Fn}{Fn}

\newtheorem{theorem}{Theorem}
\newtheorem{proposition}{Proposition}

\theoremstyle{definition}
\newtheorem{exercise}{Exercise}[subsection]
\newenvironment{exercisebreak}[1][]{%
  \begin{exercise}[#1]$ $\par\nobreak\ignorespaces
}{%
  \end{exercise}
}
\newenvironment{proofbreak}[1][\proofname]{%
  \begin{proof}[#1]$ $\par\nobreak\ignorespaces
}{%
  \end{proof}
}

\begin{document}

\title{Solutions to Exercises from Introduction To Set Theory}
\date{}
\maketitle

\tableofcontents
\newpage

\section{Sets}

\setcounter{subsection}{2}

\subsection{The Axioms}

\setcounter{exercise}{6}

\begin{exercise}
  Prove the Axiom of Pair, the Axiom of Union, and the Axiom of Power Set using
  their weaker versions.

  \begin{itemize}
  \item Weak Axiom of Pair. For any \(A\) and \(B\), there is a set \(C\) such
    that \(A \in C\) and \(B \in C\).
  \item Weak Axiom of Union. For any \(S\), there exists \(U\) such that if
    \(x \in A\) and \(A \in S\), then \(x \in U\).
  \item Weak Axiom of Power Set. For any set \(S\), there exists a \(P\) such
    that \(x \subseteq S\) implies \(x \in P\).
  \end{itemize}
\end{exercise}

\begin{proofbreak}
  \begin{itemize}
  \item The Axiom of Pair. Let \(A\) and \(B\) be sets, and let \(C\) be a set
    such that \(A \in C\) and \(B \in C\). Now consider the set
    \(Y = \left\{ x \in C : x = A \text{ or } x = B \right\}\). We show that
    \(x \in Y\) if and only if \(x = A\) or \(x = B\). Indeed, \(x = A\) or
    \(x = B\) for any \(x \in Y\). Conversely, if \(x = A\) or \(x = B\), then
    \(x \in C\) by the definition of \(C\), giving \(x \in Y\). Hence \(Y\) has the
    required property.
  \item The Axiom of Union. Let \(S\) be a set. By the weak axiom of union,
    there exists a \(U\) such that for every \(x\) and \(A\), \(x \in U\) if
    \(x \in A\) and \(A \in S\). Define
    \(Y = \left\{ x \in U : x \in A \text{ and } A \in S \text{ for some } A
    \right\}\). Clearly, if \(x \in Y\), then \(x \in A\) and \(A \in S\) for some
    \(A\). Conversely, if \(x\) is a set such that \(x \in A\) and \(A \in S\) for
    some \(A\), then \(x \in U\), by definition. Hence \(x \in Y\). Therefore, for any
    \(S\), there is a \(Y\) such that \(x \in Y\) if and only if \(x \in A\) and
    \(x \in S\) for some \(A\).
  \item The Axiom of Power Set. Let \(S\) be a set, and let \(P\) be the set
    implied by the weak axiom of power set. Consider the set
    \(Y = \left\{ x \in P : x \subseteq S \right\}\). Since \(x \subseteq S\) implies
    \(x \in P\), the statements "\(x \subseteq S\)" and
    "\(x \in P \text{ and } x \subseteq S\)" are equivalent. Hence we have
    \(x \in Y \leftrightarrow x \in P \text{ and } x \subseteq S \leftrightarrow x \subseteq S\). Therefore
    \(Y\) satisfies the required property.
  \end{itemize}
\end{proofbreak}

\subsection{Elementary Operations on Sets}

\setcounter{exercise}{3}

\begin{exercise}
  Let \(A\) be a set; show that a "complement" of \(A\) does not exist. (The
  "complement" of \(A\) is the set of all \(x \notin A\).)
\end{exercise}

\begin{proof}
  Let \(A\) be a set. Suppose that there exists a set \(X\) such that
  \(x \in X\) if and only if \(x \notin A\). Now consider the set \(A \cup X\). For every
  \(x\), \(x \in A \cup X\) if and only if \(x \in A\) or \(x \in X\), i.e., if and only
  if \(x \in A\) or \(x \notin A\). That implies \(x \in A \cup X\) for every
  \(x\) and thus \(A \cup X\) is the universal set. But this contradicts Exercise
  3.3, which states no universal set exists. Therefore the "complement" of \(A\)
  does not exist.
\end{proof}

\begin{exercise}
  Let \(S \ne \varnothing\) and \(A\) be sets.

  \begin{enumerate}
  \item Set
    \(T_1 = \left\{ Y \in \mathcal{P}(A) : Y = A \cap X \text{ for some } X \in S \right\}\), and
    prove \(A \cap \bigcup S = \bigcup T_1\) (generalized distributed law).
  \item Set \(T_2 = \left\{ Y \in \mathcal{P}(A) : Y = A - X \text{ for some } X \in S \right\}\), and prove
    \[A - \bigcup S = \bigcap T_2\]
    \[A - \bigcap S = \bigcup T_2\]
    (generalized De Morgan laws).
  \end{enumerate}
\end{exercise}

\begin{proofbreak}
  \begin{enumerate}
  \item First, we introduce a small lemma: For every \(X \in S\),
    \(A \cap X \in \mathcal{P}(A)\). Indeed, \(A \cap X \subseteq A\), and so by definition of
    \(\mathcal{P}(A)\), \(A \cap X \in \mathcal{P}(A)\).

    Now we prove that both sets are equal using the Axiom of Extensionality.

    First, let \(x \in A \cap \bigcup S\). To show that \(x \in \bigcup T_1\), we seek an element of \(T_1\)
    which contains \(x\). since \(x\) belongs to \(\bigcup S\), \(x \in X\) for some \(X \in S\).
    Now consider the set \(A \cap X\). By the above lemma, it is an element of
    \(\mathcal{P}(A)\), and thus an element of \(T_1\).
    Since \(A\) and \(X\) both contain \(x\), \(x \in A \cap X\). Therefore, \(A \cap X\)
    is the set we were seeking.

    Conversely, let \(x \in \bigcup T_1\). By the Axiom of Union, there exists a \(Y \in T_1\)
    which contains \(x\). Now by the definition of \(T_1\), \(Y = A \cap X\) for some \(X \in S\).
    Since \(x \in Y = A \cap X\), we have \(x \in \bigcup S\), as \(x \in X\) for some \(X \in S\).
    Hence \(x \in A \cap \bigcup S\), as desired.

    Therefore, by the Axiom schema of Extensionality, \(A \cap \bigcup S = \bigcup T_1\).
  \item
    \begin{enumerate}
    \item \(A - \bigcup S = \bigcap T_2\). First, suppose that
      \(x \in A - \bigcup S\). Then we have \(x \in A\), while \(x\) belongs to none of
      the sets in \(S\). Now let \(Y \in T_2\). By the definition of \(T_2\),
      there exists an \(X \in S\) such that \(Y = A - X\). Since \(x \in A\) and
      \(x \notin X\), it follows that \(x \in Y\). Hence \(x \in Y\) for every
      \(Y \in T_2\), and so \(x \in \bigcap T_2\).

      Conversely, suppose that \(x \in \bigcap T_2\). For any \(X \in S\), we have \(A - X \in T_2\)
      by definition. Then \(x \in A - X\), since \(x\) belongs to any element of \(T_2\).
      Hence for every \(X \in S\), \(x \in A\) and \(x \notin X\). Since
      there is no \(X \in S\) such that \(x \in X\), \(x \notin \bigcup S\).
      Therefore \(x \in A - \bigcup S\).

      By extensionality, \(A - \bigcup S = \bigcap T_2\).
    \item \(A - \bigcap S = \bigcup T_2\).
    \end{enumerate}
  \end{enumerate}
\end{proofbreak}

\section{Relations, Functions, and Orderings}

\subsection{Ordered Pairs}

\subsection{Relations}

\subsection{Functions}

\subsection{Equivalences and Partitions}

\subsection{Orderings}

\begin{exercisebreak}
  \begin{enumerate}
  \item Let \(R\) be an ordering of \(A\), \(S\) be the corresponding strict
    ordering of \(A\), and \(R^{*}\) be the ordering corresponding to \(S\).
    Show that \(R^{*} = R\).
  \item Let \(S\) be a strict ordering of \(A\), \(R\) be the corresponding
    ordering, and \(S^{*}\) be the strict ordering corresponding to \(R\). Then
    \(S^{*} = S\).
  \end{enumerate}
\end{exercisebreak}

\begin{exercise}
  Let \(\Fn(X, Y)\) be the set of all functions mapping a subset of \(X\) into
  \(Y\). Define a function \(\le\) in \(\Fn(X, Y)\) by

  \[f \le g \text{ if and only if } f \subseteq g.\]
  \begin{enumerate}
  \item Prove that \(\le\) is an ordering of \(\Fn(X, Y)\).
  \item Let \(F \subseteq \Fn(X, Y)\). Show that \(\sup F\) exists if and only if
    \(F\) is a compatible system of functions; then \(\sup F = \bigcup F\).
  \end{enumerate}
\end{exercise}

\begin{proofbreak}
  \begin{enumerate}
  \item The proof is obvious from the properties of \(\subseteq\).
  \item
    \begin{enumerate}
    \item Suppose that \(\text{sup}(F)\) exists, and let \(f, g \in F\). Let
      \(x \in \text{dom}(f) \cap \text{dom}(g)\). Since \(f \le \text{sup}(F)\), i.e.,
      \(f \subseteq \text{sup}(F)\), we have \(f(x) = (\text{sup}(F))(x)\). Similarly,
      \(g(x) = (\text{sup}(F))(x)\), and we conclude that \(f(x) = g(x)\).
    \item Suppose that \(F\) is a compatible system of functions. We claim that
      \(\bigcup F = \text{sup}(F)\).

      \begin{enumerate}
      \item By Theorem 3.12, \(\bigcup F\) is a function. If
        \(x \in \text{dom}(\bigcup F)\), \((x, y) \in f\) for some \(y\) and
        \(f \in F\). But then \(x \in X\), since \(\text{dom}(f) \subseteq X\). So
        \(\text{dom}(\bigcup F) \subseteq X\), and similarly,
        \(\text{ran}(\bigcup F) \subseteq Y\). Hence \(\bigcup F \in \text{Fn(X, Y)}\).
      \item Let \(f \in F\) and \((x, y) \in f\). Then immediately
        \((x, y) \in \bigcup F\), so \(f \le \bigcup F\) for every \(f \in F\);
        \(\bigcup F\) is an upper bound of \(F\).
      \item Finally, let \(g\) be an upper bound of \(F\). If
        \((x, y) \in \bigcup F\), \((x, y) \in f\) for some \(f \in F\). But then
        \((x, y) \in g\), since \(f \subseteq g\). Therefore
        \(\bigcup F \subseteq g\), that is, \(\bigcup F \le g\).
      \end{enumerate}
    \end{enumerate}
  \end{enumerate}
\end{proofbreak}
\end{document}
